\chapter{Introduction}

\section{Purpose}

Given the increase use of artificial intelligence solutions in reliably solving many of our current day-to-day problems, we decided to tackle the problem of face-based authentication in the systems that we use daily.

The challenge can be presented as follows: given a device that has access to a camera and which has been previously setup with a set of subject identities, the task is to validate the face in front of the camera as being known or not and in the positive case, the device should unlock itself.

The problems that arise are multiple and difficult to solve. Here are a few we worked on: firstly, the face in front of the camera should be validated as being real, therefore we need to give special attention to combating face spoofing attacks as printed photos, video replays or 3D masks. Secondly, the identification component of the system should be accurate enough so that two people with close facial features can be differentiated.

\section{Impact}
A successful implementation of such a solution could be deployed to border control system which would help amend the problem of crowded airports, could be used in the process of card payments as a method of reliable authentication and also could be used as an important part of an Artificial Intelligence assistant in order to create a more human-like experience.

\section{Evolution of the facial recognition systems}
The first trials of developping facial recognition systems started out in the 1960s when Woodrow Bledsoe created a semi-automated system \cite{DavisMRETO64} containing a RAND tablet which was able to uniquely describe $10^{6}$ locations in a 10x10 inch area.
An operator was required to mark specific landmarks like the individual's  hairline, eyes and nose on this tablet that were used to calculate distances and rations to a point of refference and then compared to reference data. The RAND tablet is considered to be the first graphical computer input device that is digital.

As the limits imposed by technology dropped and the potential in creating reliable facial recognition systems became worth taking into consideration, interest was shown by the Defense Advanced Research Products Agency (DARPA) which sponsored the Face Recognition Technology Evaluation (FERET) \cite{PhillipsJMHRSRP00} from 1993 to 1997 encouraging the development of face recognition algorithms and technologies by assessing the prototypes of face recognition systems. This project helped to a great extent to create a market of comercially available products.
Built upon the work of FERET, the Face Recognition Vendor Test \cite{PhillipsGMBTB03} developed by the National Institute of Standards and Technology \cite{FaceNist} and performed several times since 2000 and it provides independent government evaluations of commercially available and prototype face recognition technologies and also helps itentify future research direction for the face recognition scientific comunity.



Relative to the one of the most reliable tests in the field based on the Labeled Faces in the Wild database, the current state of the art is held by deep learning approaches. We enlist some of the best results:
\begin{itemize}
	\item PingAnLab89 0.9980 0.0016
	\item Baidu64 0.9977 0.0006
	\item FaceNet62 0.9963 0.0016
\end{itemize}



Conținut secțiune, exemplu citare \cite{hoare_csp}.
\textit{Cuvânt pentru} \textbf{glosar} \index{glosar}.	

\begin{definition}
	Conținut definiție
\end{definition}

\begin{lemma}
	Conținut lemă.
\end{lemma}

\begin{theorem}
	Conținut teoremă.
	\[
		e^{i * \pi} + 1 = 0
	\]
\end{theorem}

\begin{remark}
	Conținut observație
\end{remark}

Ecuație inline, $2^{10} = 1024 $

\begin{figure}[h]
	\begin{center}
			\includegraphics[width=4cm]{university}
	\end{center}
	\caption{Captură imagine}
\end{figure}